\chapter{Stan wiedzy}\label{ch:stan-wiedzy}

% TODO: przegląd literatury i rozwiązań — ok. 8--15 stron.

\section{Rozpoznawanie mowy (Speech-to-Text)}\label{sec:stt}

% TODO: modele Whisper, faster-whisper, wav2vec2, diaryzacja mówców,
% wyzwania języka polskiego w~STT, aktualne benchmarki WER/CER.

\section{Małe modele językowe (Small Language Models)}\label{sec:slm}

% TODO: lokalne SLM-y (Phi, Llama, Bielik), kwantyzacja,
% zastosowania w~strukturyzacji tekstu medycznego, format SOAP.

\section{Anonimizacja danych osobowych}\label{sec:anonimizacja}

% TODO: Microsoft Presidio, rozpoznawacze PII dla języka polskiego,
% PESEL/NIP/adres, wyzwania domeny medycznej.

\section{Retrieval-Augmented Generation}\label{sec:rag}

% TODO: RAG, embeddingi dla języka polskiego, bazy wektorowe (ChromaDB, Qdrant),
% hybrydowe wyszukiwanie dense + sparse.

\section{Systemy wspomagania decyzji klinicznych}\label{sec:cdss-stan}

% TODO: istniejące rozwiązania CDSS, Epic, Nuance DAX, Abridge, Suki,
% brak rozwiązań dla języka polskiego i~kontekstu refundacji NFZ.

\section{Architektury wieloagentowe}\label{sec:multi-agent}

% TODO: LangGraph, LangChain, orkiestracja agentów, wzorce projektowe.

\section{Podsumowanie}\label{sec:stan-wiedzy-podsumowanie}

% TODO: zidentyfikowana luka — brak zintegrowanego rozwiązania CDSS dla PL
% z~natywnym wsparciem NFZ i~RODO.
